\SourceLeaf{456}%
\noindent welche nach Benutzung der durch $\mathfrak{E},\mathfrak{F},\mathfrak{G}$ und $\mathfrak{E}',\mathfrak{F}',\mathfrak{G}'$ bezeichneten Werthe die nachstehende Gleichung herbeiführen:
\[
\left(\frac{\partial u}{\partial\zeta'}\right)^2=\left(\frac{\partial u'}{\partial\zeta}\right)^2.
\]

\noindent Hiernach wird der lineare Differentialausdruck
\[
ud\zeta\pm u'd\zeta'=d\varphi
\]
\noindent das totale Differential einer Function $\varphi$ der Variabeln $\zeta,\zeta'$ darstellen. Diese Eigenschaft kann jedoch für je zwei auf einander abwickelbare Flächen $S$ und $S'$ nur bestehen, wenn das untere Vorzeichen angenommen wird, da bei entgegengesetzter Annahme im Falle des Zusammenfallens von $S$ mit $S'$ das Product $2ud\zeta$ für jede Fläche ein Totaldifferential sein müsste, ein Umstand, der nur für developpable Flächen eintritt.

In der vorstehenden Gleichung ist der von \textsc{Gauss} in No.\ 2 gewählte Ausgangspunkt für die allgemeine Lösung des Problems der Deformation enthalten.

Diese Gleichung besteht für je zwei auf einander abwickelbare Flächen $S$ und $S'$, aber \textls[140]{nicht nur} für solche.

In welcher Weise \textsc{Gauss} dieselbe für diese allgemeine Lösung zu verwerthen gedachte, lässt sich aus No.\ 4 der Bemerkung nicht ohne weiteres erkennen, und bedarf diese Frage noch weiter eingehender Untersuchungen in Betreff der Eigenschaften der im besondern Fall zu wählenden Function $\varphi$.
\begin{flushright}\textsc{Weingarten}.\end{flushright}

